% ==================================================================
% MODELO — Princípios: modelo mais simples que gera o insight;
% ambiente em texto corrido ANTES da matemática; suposições numeradas
% com conteúdo econômico explícito; proposições com intuição ANTES da
% demonstração; provas no apêndice. Cada proposição gera uma previsão
% testável mapeada a uma regressão da Seção 4.
% ==================================================================
\section{Modelo}
\label{sec:modelo}

\subsection{Ambiente}

Apresento um modelo estático de designação ocupacional em equilíbrio parcial com dois setores contratuais --- formal ($F$) e informal ($I$) --- e um contínuo de trabalhadores com escolaridade heterogênea $e \in [e_{\min}, e_{\max}]$, distribuídos segundo uma função de densidade contínua $f(e)$. A economia dispõe de $J$ ocupações distintas, indexadas por $j \in \{1, \dots, J\}$. Cada ocupação é caracterizada por um vetor fixo de tarefas com índice de exposição tecnológica à inteligência artificial $\theta_j \in [0, 10]$, determinado exogenamente pelas características intrínsecas das rotinas informacionais \citep{eloundou2024gpts}. 

Os trabalhadores escolhem simultaneamente a ocupação $j$ e o arranjo contratual $s \in \{F, I\}$ para maximizar sua remuneração líquida, tomando os parâmetros tecnológicos e os custos regulatórios como dados. Firmas do setor formal incorrem em um custo fixo de conformidade tributária e institucional $c > 0$ por trabalhador \citep{ulyssea2018firms}, mas fornecem a infraestrutura corporativa, o capital de dados e a governança que potencializam os ganhos de produtividade proporcionados pela IA. No setor informal, os trabalhadores acessam livremente ferramentas individuais de IA, mas não dispõem do capital complementar organizacional necessário para converter a tecnologia em ganhos corporativos de escala. Em termos econômicos, a informalidade amortece a exposição à automação ao limitar o \emph{retorno econômico} da integração tecnológica corporativa --- e não por barrar o acesso à ferramenta.

\subsection{Primitivos de produção e salários}

Um trabalhador com escolaridade $e$ alocado à ocupação $j$ sob o arranjo contratual $s \in \{F, I\}$ gera produto
\begin{equation}
    y_{j s}(e) = g(e) \, \big(1 + \kappa_s \, \theta_j\big),
    \label{eq:produto}
\end{equation}
onde $g(e)$ (com $g' > 0$ e $g'' \le 0$) representa o componente básico de capital humano, e $\kappa_s > 0$ parametriza a eficácia da IA sob o arranjo $s$. Defino a estrutura de retorno tecnológico por:
\begin{equation}
    \kappa_s = \kappa_{\text{ind}} + \Delta\kappa \, \mathbb{1}[s = F],
    \quad \text{com } \Delta\kappa \equiv \kappa_{\text{org}} - \kappa_{\text{ind}} > 0,
\end{equation}
em que $\kappa_{\text{ind}} > 0$ captura o ganho da IA individual autônoma (disponível universalmente) e $\Delta\kappa > 0$ reflete o \emph{prêmio de capital organizacional} restrito às empresas do setor formal.

Sob concorrência perfeita e diferenciais salariais hedônicos que equilibram oferta e demanda entre ocupações \citep{rosen1986theory,costinot2010matching}, a remuneração de equilíbrio dentro de cada ocupação iguala a produtividade líquida do custo de conformidade:
\begin{equation}
    w_{jF}(e) = g(e)\big(1 + \kappa_{\text{org}}\,\theta_j\big) - c,
    \qquad w_{jI}(e) = g(e)\big(1 + \kappa_{\text{ind}}\,\theta_j\big).
    \label{eq:salarios}
\end{equation}

\noindent O modelo fundamenta-se em três hipóteses centrais:
\begin{itemize}
    \item \textbf{H1 (Exogeneidade tecnológica de $\theta_j$).} O escore de exposição $\theta_j$ reflete atributos intrínsecos ao conteúdo de tarefas da ocupação \citep{eloundou2024gpts}, predeterminado em relação aos choques locais do mercado de trabalho brasileiro.
    \item \textbf{H2 (Prêmio de capital organizacional $\Delta\kappa > 0$).} O retorno produtivo da IA sob integração corporativa supera o uso individual informal ($\kappa_{\text{org}} > \kappa_{\text{ind}}$).
    \item \textbf{H3 (Custo de conformidade institucional $c > 0$).} A formalização de postos de trabalho impõe encargos regulatórios e parafiscais por trabalhador \citep{ulyssea2018firms}.
\end{itemize}

\subsection{Equilíbrio de designação e previsões testáveis}

Um \emph{equilíbrio de designação} em equilíbrio parcial é definido por uma função de alocação $\sigma: [e_{\min}, e_{\max}] \to \{1, \dots, J\} \times \{F, I\}$ e um envelope salarial $w^*(e) = \max_{j, s} w_{js}(e)$ tais que cada trabalhador escolhe a ocupação e o regime contratual ótimos.

Dentro de qualquer ocupação $j$, a escolha entre formalidade e informalidade obedece à regra de arbitragem líquida:
\begin{equation}
    w_{jF}(e) \ge w_{jI}(e) \iff g(e) \, \Delta\kappa \, \theta_j \ \ge \ c.
    \label{eq:corte}
\end{equation}
Sendo $g(\cdot)$ estritamente crescente e invertível, a condição \eqref{eq:corte} determina um limiar único de escolaridade para formalização na ocupação $j$:
\begin{equation}
    e^*_j = g^{-1}\!\left(\frac{c}{\Delta\kappa \, \theta_j}\right),
    \label{eq:limiar}
\end{equation}
onde trabalhadores com $e \ge e^*_j$ inserem-se no setor formal e aqueles com $e < e^*_j$ operam na informalidade.

\begin{proposicao}[Gradiente escolaridade--exposição]
\label{prop:gradiente}
No equilíbrio de designação, a escolaridade média dos trabalhadores é estritamente crescente na exposição da ocupação: $\partial \bar{e}_j / \partial \theta_j > 0$.
\end{proposicao}

\textit{Intuição econômica.} A remuneração de reserva $w^*(e, j) \equiv \max_{s} w_{js}(e)$ é estritamente supermodular em $(e, \theta_j)$ e o limiar de formalização $e_j^*$ é decrescente em $\theta_j$; pelo pareamento seletivo positivo de \citet{becker1973theory}, trabalhadores com maior capital humano concentram-se no topo da distribuição de $\theta_j$. \emph{Previsão empírica:} coeficiente $\beta_1 > 0$ na regressão de exposição sobre escolaridade (especificação S1).

\begin{proposicao}[Complementaridade exposição--formalidade]
\label{prop:formalidade}
A taxa de formalidade dentro da ocupação é estritamente crescente na exposição: $\partial \, \text{formal}_j / \partial \theta_j > 0$.
\end{proposicao}

\textit{Intuição econômica.} Pela equação \eqref{eq:limiar}, $e^*_j$ decresce estritamente em $\theta_j$: em ocupações altamente expostas, o prêmio de produtividade organizacional compensa o custo regulatório $c$ para uma faixa mais ampla de escolaridade, expandindo a massa de postos formais. \emph{Previsão empírica:} coeficiente negativo e significante da exposição sobre a probabilidade de informalidade na relação sem escolaridade (especificação S3a).

\begin{corolario}[Mediação parcial pela composição educacional]
\label{cor:mediacao}
A associação negativa bruta entre exposição à IA e probabilidade de informalidade é substancialmente atenuada quando se condiciona na escolaridade individual do trabalhador, retendo uma associação residual autônoma decorrente de complementaridades organizacionais de tarefas.
\end{corolario}

\textit{Intuição econômica.} Como $e_j^*$ é decrescente em $\theta_j$, ocupações mais expostas atraem trabalhadores mais escolarizados (Proposição~\ref{prop:gradiente}), de modo que a composição educacional responde por parcela substantiva da menor informalidade observada nelas. Como a integração da tecnologia requer ativos corporativos específicos ($\Delta\kappa\,\theta_j$), a exposição mantém um canal autônomo de formalização mesmo condicionando na escolaridade individual. \emph{Previsão empírica:} em S3, o coeficiente de exposição atenua-se expressivamente em relação a S3a, mas permanece negativo e significante ($\delta_1 < 0$).

\subsection{Discussão teórica}

Em suma, a difusão de ferramentas de IA generativa não anula a função amortecedora da informalidade: a reorganização produtiva em larga escala depende de ativos organizacionais complementares cuja viabilidade econômica restringe-se ao ambiente formal.
